\begin{abstract}

    \begin{center}
        \large{% Единое место для указания названия работы
Влияние графовой структуры на способность больших языковых моделей вести цепочку рассуждений
} \\
    \large\textit{Орлов Алексей Алексеевич} \\[1 cm]
    \end{center}

    В работе исследуется влияние явной графовой структуры на качество многошаговых рассуждений больших языковых моделей.
    Для автоматизированной оценки структур ответов разработан программный комплекс 
    \mbox{smart-thinking-llm},
    обеспечивающий преобразование текстовых генераций в графовое представление, их сопоставление с эталонными путями и 
    расчёт метрик покрытия на датасетах RuleTaker и Wikidata. 
    Проведён сравнительный анализ рассуждающих способностей 
    ряда моделей (Qwen3, Llama 3.x, Olmo-Think, GPT-OSS), по результатам которого для дообучения отобрана архитектура Qwen3-0.6B. 
    С использованием платформы TunePPO реализован двухэтапный пайплайн обучения, в ходе которого исследовано влияние 
    специальных функций награды, регулирующих структурное соответствие графа и корректность финального ответа. 
    Полученные экспериментальные результаты демонстрируют, что интеграция графовой 
    награды статистически значимо повышает точность логических выводов и снижает долю структурных разрывов по сравнению 
    с базовым подходом, опирающимся исключительно на бинарную оценку результата. Полученные данные подтверждают эффективность 
    структурно-ориентированного выравнивания компактных LLM и предлагают воспроизводимую методологию оценки качества цепочек 
    рассуждений.

    \vfill
    % \begin{center}
    % \textbf{Abstract} \\[1 cm]

    % Research and development of machine learning methods
    % \end{center}

\end{abstract}
